%%%%%%%%%%%%%%%%%%%%%%%%%%%%%%%%%%%%%%%%%%%%%%%%%%%%%%%%%%%%
% 14.5 ALGEBRAIC IDENTITIES
% The Composition of Advanced Mathematics
%%%%%%%%%%%%%%%%%%%%%%%%%%%%%%%%%%%%%%%%%%%%%%%%%%%%%%%%%%%%

\section{Algebraic Identities}
\label{sec:algebraic-identities}

\prerequisite{
Arithmetic, algebraic notation, exponents, radicals, logarithms,
polynomials, rational expressions, and basic complex numbers.
}

\learningobjectives{
By the end of this reference section, the reader should be able to:
\begin{itemize}
    \item recognize the distinction between an identity and an equation;
    \item use exponent laws correctly;
    \item simplify powers and radicals;
    \item apply logarithmic identities;
    \item recognize common polynomial identities;
    \item expand binomial expressions;
    \item factor differences of squares;
    \item factor sums and differences of cubes;
    \item use perfect-square and perfect-cube identities;
    \item recognize generalized difference-of-powers identities;
    \item apply the binomial theorem;
    \item use multinomial expansions conceptually;
    \item manipulate rational expressions;
    \item simplify nested fractions;
    \item rationalize denominators when appropriate;
    \item use absolute-value identities correctly;
    \item recognize identities involving complex conjugates;
    \item distinguish valid identities from common false identities;
    \item identify domain restrictions attached to algebraic transformations;
    \item use algebraic identities efficiently in calculus, analysis,
          number theory, and applied mathematics.
\end{itemize}
}

%%%%%%%%%%%%%%%%%%%%%%%%%%%%%%%%%%%%%%%%%%%%%%%%%%%%%%%%%%%%
\subsection{What Is an Algebraic Identity?}
\label{subsec:what-is-algebraic-identity}
%%%%%%%%%%%%%%%%%%%%%%%%%%%%%%%%%%%%%%%%%%%%%%%%%%%%%%%%%%%%

An algebraic identity is an equality that is true for every value of the
variables for which both sides are defined.

For example,

\[
\boxed{
(a+b)^2
=
a^2+2ab+b^2
}
\]

is an identity.

By contrast,

\[
x^2=4
\]

is an equation that is true only for selected values of \(x\).

%%%%%%%%%%%%%%%%%%%%%%%%%%%%%%%%%%%%%%%%%%%%%%%%%%%%%%%%%%%%
\subsection{Identity Versus Equation}
\label{subsec:identity-versus-equation}
%%%%%%%%%%%%%%%%%%%%%%%%%%%%%%%%%%%%%%%%%%%%%%%%%%%%%%%%%%%%

An identity may be written

\[
\boxed{
L(x)=R(x)
\quad
\text{for all admissible }x.
}
\]

An equation asks for values satisfying

\[
\boxed{
L(x)=R(x).
}
\]

For example,

\[
\boxed{
(x+1)^2
=
x^2+2x+1
}
\]

is an identity, while

\[
\boxed{
(x+1)^2=9
}
\]

is an equation.

%%%%%%%%%%%%%%%%%%%%%%%%%%%%%%%%%%%%%%%%%%%%%%%%%%%%%%%%%%%%
\subsection{Commutative Laws}
\label{subsec:commutative-laws-reference}
%%%%%%%%%%%%%%%%%%%%%%%%%%%%%%%%%%%%%%%%%%%%%%%%%%%%%%%%%%%%

For ordinary real or complex numbers,

\[
\boxed{
a+b=b+a
}
\]

and

\[
\boxed{
ab=ba.
}
\]

These are the commutative laws for addition and multiplication.

They do not automatically hold for every algebraic structure.

For example, matrix multiplication is generally not commutative.

%%%%%%%%%%%%%%%%%%%%%%%%%%%%%%%%%%%%%%%%%%%%%%%%%%%%%%%%%%%%
\subsection{Associative Laws}
\label{subsec:associative-laws-reference}
%%%%%%%%%%%%%%%%%%%%%%%%%%%%%%%%%%%%%%%%%%%%%%%%%%%%%%%%%%%%

For ordinary numbers,

\[
\boxed{
(a+b)+c
=
a+(b+c)
}
\]

and

\[
\boxed{
(ab)c
=
a(bc).
}
\]

These are the associative laws.

They justify omitting unnecessary parentheses in repeated addition and
multiplication.

%%%%%%%%%%%%%%%%%%%%%%%%%%%%%%%%%%%%%%%%%%%%%%%%%%%%%%%%%%%%
\subsection{Distributive Law}
\label{subsec:distributive-law-reference}
%%%%%%%%%%%%%%%%%%%%%%%%%%%%%%%%%%%%%%%%%%%%%%%%%%%%%%%%%%%%

The distributive law is

\[
\boxed{
a(b+c)
=
ab+ac.
}
\]

Similarly,

\[
\boxed{
(a+b)c
=
ac+bc.
}
\]

This identity is the foundation of polynomial expansion and factoring.

%%%%%%%%%%%%%%%%%%%%%%%%%%%%%%%%%%%%%%%%%%%%%%%%%%%%%%%%%%%%
\subsection{Negation Rules}
\label{subsec:negation-rules-reference}
%%%%%%%%%%%%%%%%%%%%%%%%%%%%%%%%%%%%%%%%%%%%%%%%%%%%%%%%%%%%

Important identities include

\[
\boxed{
-(a+b)
=
-a-b,
}
\]

\[
\boxed{
-(a-b)
=
-a+b,
}
\]

and

\[
\boxed{
(-a)(-b)=ab.
}
\]

Also,

\[
\boxed{
(-a)b
=
-(ab).
}
\]

%%%%%%%%%%%%%%%%%%%%%%%%%%%%%%%%%%%%%%%%%%%%%%%%%%%%%%%%%%%%
\subsection{Zero and One}
\label{subsec:zero-one-identities}
%%%%%%%%%%%%%%%%%%%%%%%%%%%%%%%%%%%%%%%%%%%%%%%%%%%%%%%%%%%%

For any number \(a\),

\[
\boxed{
a+0=a
}
\]

and

\[
\boxed{
a\cdot1=a.
}
\]

Also,

\[
\boxed{
a-a=0
}
\]

and, when \(a\neq0\),

\[
\boxed{
\frac{a}{a}=1.
}
\]

%%%%%%%%%%%%%%%%%%%%%%%%%%%%%%%%%%%%%%%%%%%%%%%%%%%%%%%%%%%%
\subsection{Exponent Laws}
\label{subsec:exponent-laws-reference}
%%%%%%%%%%%%%%%%%%%%%%%%%%%%%%%%%%%%%%%%%%%%%%%%%%%%%%%%%%%%

For appropriate bases and integer exponents,

\[
\boxed{
a^ma^n
=
a^{m+n}.
}
\]

Also,

\[
\boxed{
\frac{a^m}{a^n}
=
a^{m-n},
\qquad
a\neq0,
}
\]

and

\[
\boxed{
(a^m)^n
=
a^{mn}.
}
\]

%%%%%%%%%%%%%%%%%%%%%%%%%%%%%%%%%%%%%%%%%%%%%%%%%%%%%%%%%%%%
\subsection{Product and Quotient Powers}
\label{subsec:product-quotient-powers}
%%%%%%%%%%%%%%%%%%%%%%%%%%%%%%%%%%%%%%%%%%%%%%%%%%%%%%%%%%%%

For integer \(n\),

\[
\boxed{
(ab)^n
=
a^nb^n.
}
\]

If \(b\neq0\),

\[
\boxed{
\left(
\frac{a}{b}
\right)^n
=
\frac{a^n}{b^n}.
}
\]

%%%%%%%%%%%%%%%%%%%%%%%%%%%%%%%%%%%%%%%%%%%%%%%%%%%%%%%%%%%%
\subsection{Zero Exponent}
\label{subsec:zero-exponent-reference}
%%%%%%%%%%%%%%%%%%%%%%%%%%%%%%%%%%%%%%%%%%%%%%%%%%%%%%%%%%%%

For

\[
a\neq0,
\]

\[
\boxed{
a^0=1.
}
\]

This follows from

\[
\frac{a^m}{a^m}
=
a^{m-m}
=
a^0
\]

and

\[
\frac{a^m}{a^m}=1.
\]

%%%%%%%%%%%%%%%%%%%%%%%%%%%%%%%%%%%%%%%%%%%%%%%%%%%%%%%%%%%%
\subsection{Negative Exponents}
\label{subsec:negative-exponents-reference}
%%%%%%%%%%%%%%%%%%%%%%%%%%%%%%%%%%%%%%%%%%%%%%%%%%%%%%%%%%%%

For \(a\neq0\),

\[
\boxed{
a^{-n}
=
\frac{1}{a^n}.
}
\]

Thus

\[
\boxed{
\frac{1}{a^{-n}}
=
a^n.
}
\]

%%%%%%%%%%%%%%%%%%%%%%%%%%%%%%%%%%%%%%%%%%%%%%%%%%%%%%%%%%%%
\subsection{Rational Exponents}
\label{subsec:rational-exponents-reference}
%%%%%%%%%%%%%%%%%%%%%%%%%%%%%%%%%%%%%%%%%%%%%%%%%%%%%%%%%%%%

For suitable real \(a\),

\[
\boxed{
a^{1/n}
=
\sqrt[n]{a}.
}
\]

More generally,

\[
\boxed{
a^{m/n}
=
\sqrt[n]{a^m}
=
\left(
\sqrt[n]{a}
\right)^m.
}
\]

When \(n\) is even and one works over the real numbers, domain restrictions
must be observed.

%%%%%%%%%%%%%%%%%%%%%%%%%%%%%%%%%%%%%%%%%%%%%%%%%%%%%%%%%%%%
\subsection{Worked Example: Exponent Simplification}
\label{subsec:worked-exponent-simplification}
%%%%%%%%%%%%%%%%%%%%%%%%%%%%%%%%%%%%%%%%%%%%%%%%%%%%%%%%%%%%

\begin{workedexamplebox}{Simplify an Exponential Expression}

Simplify

\[
\frac{x^5x^{-2}}{x}.
\]

Combine powers in the numerator:

\[
x^5x^{-2}
=
x^3.
\]

Then

\[
\frac{x^3}{x}
=
x^2,
\qquad
x\neq0.
\]

\begin{answerbox}
\[
\boxed{
\frac{x^5x^{-2}}{x}
=
x^2,
\qquad
x\neq0.
}
\]
\end{answerbox}

\end{workedexamplebox}

%%%%%%%%%%%%%%%%%%%%%%%%%%%%%%%%%%%%%%%%%%%%%%%%%%%%%%%%%%%%
\subsection{Radical Product Rule}
\label{subsec:radical-product-rule}
%%%%%%%%%%%%%%%%%%%%%%%%%%%%%%%%%%%%%%%%%%%%%%%%%%%%%%%%%%%%

For nonnegative real numbers \(a,b\),

\[
\boxed{
\sqrt{ab}
=
\sqrt{a}\sqrt{b}.
}
\]

More generally, for appropriate domains,

\[
\boxed{
\sqrt[n]{ab}
=
\sqrt[n]{a}
\sqrt[n]{b}.
}
\]

%%%%%%%%%%%%%%%%%%%%%%%%%%%%%%%%%%%%%%%%%%%%%%%%%%%%%%%%%%%%
\subsection{Radical Quotient Rule}
\label{subsec:radical-quotient-rule}
%%%%%%%%%%%%%%%%%%%%%%%%%%%%%%%%%%%%%%%%%%%%%%%%%%%%%%%%%%%%

For

\[
a\geq0,
\qquad
b>0,
\]

\[
\boxed{
\sqrt{
\frac{a}{b}
}
=
\frac{\sqrt a}{\sqrt b}.
}
\]

%%%%%%%%%%%%%%%%%%%%%%%%%%%%%%%%%%%%%%%%%%%%%%%%%%%%%%%%%%%%
\subsection{Square Root of a Square}
\label{subsec:square-root-square}
%%%%%%%%%%%%%%%%%%%%%%%%%%%%%%%%%%%%%%%%%%%%%%%%%%%%%%%%%%%%

Over the real numbers,

\[
\boxed{
\sqrt{x^2}
=
|x|.
}
\]

It is not generally true that

\[
\sqrt{x^2}=x.
\]

For example,

\[
\sqrt{(-3)^2}
=
3,
\]

not \(-3\).

%%%%%%%%%%%%%%%%%%%%%%%%%%%%%%%%%%%%%%%%%%%%%%%%%%%%%%%%%%%%
\subsection{Even and Odd Roots}
\label{subsec:even-odd-roots}
%%%%%%%%%%%%%%%%%%%%%%%%%%%%%%%%%%%%%%%%%%%%%%%%%%%%%%%%%%%%

For odd positive integer \(n\),

\[
\boxed{
\sqrt[n]{x^n}=x
}
\]

for every real \(x\).

For even positive integer \(n\),

\[
\boxed{
\sqrt[n]{x^n}
=
|x|.
}
\]

%%%%%%%%%%%%%%%%%%%%%%%%%%%%%%%%%%%%%%%%%%%%%%%%%%%%%%%%%%%%
\subsection{Combining Like Radicals}
\label{subsec:combining-like-radicals}
%%%%%%%%%%%%%%%%%%%%%%%%%%%%%%%%%%%%%%%%%%%%%%%%%%%%%%%%%%%%

Radicals can be combined only when their radical parts agree.

For example,

\[
\boxed{
3\sqrt2+5\sqrt2
=
8\sqrt2.
}
\]

But

\[
\boxed{
\sqrt2+\sqrt3
}
\]

does not simplify by combining the terms.

%%%%%%%%%%%%%%%%%%%%%%%%%%%%%%%%%%%%%%%%%%%%%%%%%%%%%%%%%%%%
\subsection{Rationalizing a Simple Denominator}
\label{subsec:rationalizing-simple-denominator}
%%%%%%%%%%%%%%%%%%%%%%%%%%%%%%%%%%%%%%%%%%%%%%%%%%%%%%%%%%%%

For

\[
a\neq0,
\]

\[
\frac{1}{\sqrt a}
\]

may be rationalized by multiplying numerator and denominator by \(\sqrt a\):

\[
\boxed{
\frac{1}{\sqrt a}
=
\frac{\sqrt a}{a}.
}
\]

%%%%%%%%%%%%%%%%%%%%%%%%%%%%%%%%%%%%%%%%%%%%%%%%%%%%%%%%%%%%
\subsection{Conjugate Rationalization}
\label{subsec:conjugate-rationalization}
%%%%%%%%%%%%%%%%%%%%%%%%%%%%%%%%%%%%%%%%%%%%%%%%%%%%%%%%%%%%

For expressions involving

\[
a+\sqrt b,
\]

the conjugate is

\[
\boxed{
a-\sqrt b.
}
\]

Since

\[
(a+\sqrt b)(a-\sqrt b)
=
a^2-b,
\]

the conjugate can remove radicals from denominators.

%%%%%%%%%%%%%%%%%%%%%%%%%%%%%%%%%%%%%%%%%%%%%%%%%%%%%%%%%%%%
\subsection{Worked Example: Rationalizing a Denominator}
\label{subsec:worked-rationalize-denominator}
%%%%%%%%%%%%%%%%%%%%%%%%%%%%%%%%%%%%%%%%%%%%%%%%%%%%%%%%%%%%

\begin{workedexamplebox}{Use a Conjugate}

Simplify

\[
\frac{1}{3+\sqrt2}.
\]

Multiply by

\[
\frac{3-\sqrt2}{3-\sqrt2}.
\]

Then

\[
\frac{1}{3+\sqrt2}
=
\frac{3-\sqrt2}{9-2}.
\]

Therefore,

\begin{answerbox}
\[
\boxed{
\frac{1}{3+\sqrt2}
=
\frac{3-\sqrt2}{7}.
}
\]
\end{answerbox}

\end{workedexamplebox}

%%%%%%%%%%%%%%%%%%%%%%%%%%%%%%%%%%%%%%%%%%%%%%%%%%%%%%%%%%%%
\subsection{Logarithm Definition}
\label{subsec:logarithm-definition-reference}
%%%%%%%%%%%%%%%%%%%%%%%%%%%%%%%%%%%%%%%%%%%%%%%%%%%%%%%%%%%%

For

\[
b>0,
\qquad
b\neq1,
\]

\[
\boxed{
\log_b x=y
\iff
b^y=x.
}
\]

The logarithm is defined only for

\[
\boxed{
x>0
}
\]

when working over the real numbers.

%%%%%%%%%%%%%%%%%%%%%%%%%%%%%%%%%%%%%%%%%%%%%%%%%%%%%%%%%%%%
\subsection{Logarithm Product Rule}
\label{subsec:logarithm-product-rule}
%%%%%%%%%%%%%%%%%%%%%%%%%%%%%%%%%%%%%%%%%%%%%%%%%%%%%%%%%%%%

For positive \(x,y\),

\[
\boxed{
\log_b(xy)
=
\log_bx+\log_by.
}
\]

For the natural logarithm,

\[
\boxed{
\ln(xy)
=
\ln x+\ln y.
}
\]

%%%%%%%%%%%%%%%%%%%%%%%%%%%%%%%%%%%%%%%%%%%%%%%%%%%%%%%%%%%%
\subsection{Logarithm Quotient Rule}
\label{subsec:logarithm-quotient-rule}
%%%%%%%%%%%%%%%%%%%%%%%%%%%%%%%%%%%%%%%%%%%%%%%%%%%%%%%%%%%%

For positive \(x,y\),

\[
\boxed{
\log_b
\left(
\frac{x}{y}
\right)
=
\log_bx-\log_by.
}
\]

%%%%%%%%%%%%%%%%%%%%%%%%%%%%%%%%%%%%%%%%%%%%%%%%%%%%%%%%%%%%
\subsection{Logarithm Power Rule}
\label{subsec:logarithm-power-rule}
%%%%%%%%%%%%%%%%%%%%%%%%%%%%%%%%%%%%%%%%%%%%%%%%%%%%%%%%%%%%

For appropriate \(x>0\),

\[
\boxed{
\log_b(x^r)
=
r\log_bx.
}
\]

In particular,

\[
\boxed{
\ln(x^r)
=
r\ln x.
}
\]

%%%%%%%%%%%%%%%%%%%%%%%%%%%%%%%%%%%%%%%%%%%%%%%%%%%%%%%%%%%%
\subsection{Change of Base}
\label{subsec:change-of-base-reference}
%%%%%%%%%%%%%%%%%%%%%%%%%%%%%%%%%%%%%%%%%%%%%%%%%%%%%%%%%%%%

For valid bases \(a,b\),

\[
\boxed{
\log_bx
=
\frac{\log_ax}{\log_ab}.
}
\]

Using natural logarithms,

\[
\boxed{
\log_bx
=
\frac{\ln x}{\ln b}.
}
\]

%%%%%%%%%%%%%%%%%%%%%%%%%%%%%%%%%%%%%%%%%%%%%%%%%%%%%%%%%%%%
\subsection{Inverse Exponential-Logarithm Identities}
\label{subsec:inverse-log-exp-identities}
%%%%%%%%%%%%%%%%%%%%%%%%%%%%%%%%%%%%%%%%%%%%%%%%%%%%%%%%%%%%

For suitable domains,

\[
\boxed{
b^{\log_bx}=x
}
\]

and

\[
\boxed{
\log_b(b^x)=x.
}
\]

For natural logs,

\[
\boxed{
e^{\ln x}=x,
\qquad
x>0,
}
\]

and

\[
\boxed{
\ln(e^x)=x.
}
\]

%%%%%%%%%%%%%%%%%%%%%%%%%%%%%%%%%%%%%%%%%%%%%%%%%%%%%%%%%%%%
\subsection{Common False Logarithm Identities}
\label{subsec:false-log-identities}
%%%%%%%%%%%%%%%%%%%%%%%%%%%%%%%%%%%%%%%%%%%%%%%%%%%%%%%%%%%%

In general,

\[
\boxed{
\ln(x+y)
\neq
\ln x+\ln y.
}
\]

Also,

\[
\boxed{
\ln(x-y)
\neq
\ln x-\ln y.
}
\]

Logarithm rules apply to products, quotients, and powers, not sums and
differences.

%%%%%%%%%%%%%%%%%%%%%%%%%%%%%%%%%%%%%%%%%%%%%%%%%%%%%%%%%%%%
\subsection{Square of a Sum}
\label{subsec:square-of-sum}
%%%%%%%%%%%%%%%%%%%%%%%%%%%%%%%%%%%%%%%%%%%%%%%%%%%%%%%%%%%%

The fundamental identity is

\[
\boxed{
(a+b)^2
=
a^2+2ab+b^2.
}
\]

This follows from

\[
(a+b)(a+b).
\]

%%%%%%%%%%%%%%%%%%%%%%%%%%%%%%%%%%%%%%%%%%%%%%%%%%%%%%%%%%%%
\subsection{Square of a Difference}
\label{subsec:square-of-difference}
%%%%%%%%%%%%%%%%%%%%%%%%%%%%%%%%%%%%%%%%%%%%%%%%%%%%%%%%%%%%

Similarly,

\[
\boxed{
(a-b)^2
=
a^2-2ab+b^2.
}
\]

The middle term is negative because

\[
a(-b)+(-b)a=-2ab.
\]

%%%%%%%%%%%%%%%%%%%%%%%%%%%%%%%%%%%%%%%%%%%%%%%%%%%%%%%%%%%%
\subsection{Difference of Squares}
\label{subsec:difference-of-squares}
%%%%%%%%%%%%%%%%%%%%%%%%%%%%%%%%%%%%%%%%%%%%%%%%%%%%%%%%%%%%

One of the most useful factoring identities is

\[
\boxed{
a^2-b^2
=
(a-b)(a+b).
}
\]

This identity works in both directions:

\[
\boxed{
(a-b)(a+b)
=
a^2-b^2.
}
\]

%%%%%%%%%%%%%%%%%%%%%%%%%%%%%%%%%%%%%%%%%%%%%%%%%%%%%%%%%%%%
\subsection{Worked Example: Difference of Squares}
\label{subsec:worked-difference-squares}
%%%%%%%%%%%%%%%%%%%%%%%%%%%%%%%%%%%%%%%%%%%%%%%%%%%%%%%%%%%%

\begin{workedexamplebox}{Factor a Polynomial}

Factor

\[
9x^2-25.
\]

Recognize

\[
9x^2=(3x)^2
\]

and

\[
25=5^2.
\]

Thus

\begin{answerbox}
\[
\boxed{
9x^2-25
=
(3x-5)(3x+5).
}
\]
\end{answerbox}

\end{workedexamplebox}

%%%%%%%%%%%%%%%%%%%%%%%%%%%%%%%%%%%%%%%%%%%%%%%%%%%%%%%%%%%%
\subsection{Cube of a Sum}
\label{subsec:cube-of-sum}
%%%%%%%%%%%%%%%%%%%%%%%%%%%%%%%%%%%%%%%%%%%%%%%%%%%%%%%%%%%%

\[
\boxed{
(a+b)^3
=
a^3+3a^2b+3ab^2+b^3.
}
\]

The coefficients

\[
1,3,3,1
\]

come from the third row of Pascal's triangle.

%%%%%%%%%%%%%%%%%%%%%%%%%%%%%%%%%%%%%%%%%%%%%%%%%%%%%%%%%%%%
\subsection{Cube of a Difference}
\label{subsec:cube-of-difference}
%%%%%%%%%%%%%%%%%%%%%%%%%%%%%%%%%%%%%%%%%%%%%%%%%%%%%%%%%%%%

\[
\boxed{
(a-b)^3
=
a^3-3a^2b+3ab^2-b^3.
}
\]

%%%%%%%%%%%%%%%%%%%%%%%%%%%%%%%%%%%%%%%%%%%%%%%%%%%%%%%%%%%%
\subsection{Sum of Cubes}
\label{subsec:sum-of-cubes}
%%%%%%%%%%%%%%%%%%%%%%%%%%%%%%%%%%%%%%%%%%%%%%%%%%%%%%%%%%%%

\[
\boxed{
a^3+b^3
=
(a+b)
(a^2-ab+b^2).
}
\]

%%%%%%%%%%%%%%%%%%%%%%%%%%%%%%%%%%%%%%%%%%%%%%%%%%%%%%%%%%%%
\subsection{Difference of Cubes}
\label{subsec:difference-of-cubes}
%%%%%%%%%%%%%%%%%%%%%%%%%%%%%%%%%%%%%%%%%%%%%%%%%%%%%%%%%%%%

\[
\boxed{
a^3-b^3
=
(a-b)
(a^2+ab+b^2).
}
\]

%%%%%%%%%%%%%%%%%%%%%%%%%%%%%%%%%%%%%%%%%%%%%%%%%%%%%%%%%%%%
\subsection{Worked Example: Sum of Cubes}
\label{subsec:worked-sum-cubes}
%%%%%%%%%%%%%%%%%%%%%%%%%%%%%%%%%%%%%%%%%%%%%%%%%%%%%%%%%%%%

\begin{workedexamplebox}{Factor \(x^3+8\)}

Since

\[
8=2^3,
\]

we have

\[
x^3+8
=
x^3+2^3.
\]

Use the sum-of-cubes identity:

\[
x^3+2^3
=
(x+2)(x^2-2x+4).
\]

\begin{answerbox}
\[
\boxed{
x^3+8
=
(x+2)(x^2-2x+4).
}
\]
\end{answerbox}

\end{workedexamplebox}

%%%%%%%%%%%%%%%%%%%%%%%%%%%%%%%%%%%%%%%%%%%%%%%%%%%%%%%%%%%%
\subsection{General Difference of Powers}
\label{subsec:general-difference-powers}
%%%%%%%%%%%%%%%%%%%%%%%%%%%%%%%%%%%%%%%%%%%%%%%%%%%%%%%%%%%%

For positive integer \(n\),

\[
\boxed{
a^n-b^n
=
(a-b)
\left(
a^{n-1}
+
a^{n-2}b
+
\cdots
+
ab^{n-2}
+
b^{n-1}
\right).
}
\]

Thus

\[
\boxed{
a-b
}
\]

always divides

\[
\boxed{
a^n-b^n.
}
\]

%%%%%%%%%%%%%%%%%%%%%%%%%%%%%%%%%%%%%%%%%%%%%%%%%%%%%%%%%%%%
\subsection{General Sum of Odd Powers}
\label{subsec:general-sum-odd-powers}
%%%%%%%%%%%%%%%%%%%%%%%%%%%%%%%%%%%%%%%%%%%%%%%%%%%%%%%%%%%%

If \(n\) is odd,

\[
\boxed{
a^n+b^n
}
\]

is divisible by

\[
\boxed{
a+b.
}
\]

More explicitly,

\[
\boxed{
a^n+b^n
=
(a+b)
\left(
a^{n-1}
-a^{n-2}b
+a^{n-3}b^2
-\cdots
-ab^{n-2}
+b^{n-1}
\right).
}
\]

This factorization does not hold in the same form for even \(n\).

%%%%%%%%%%%%%%%%%%%%%%%%%%%%%%%%%%%%%%%%%%%%%%%%%%%%%%%%%%%%
\subsection{Perfect-Square Trinomials}
\label{subsec:perfect-square-trinomials}
%%%%%%%%%%%%%%%%%%%%%%%%%%%%%%%%%%%%%%%%%%%%%%%%%%%%%%%%%%%%

The patterns

\[
\boxed{
a^2+2ab+b^2
=
(a+b)^2
}
\]

and

\[
\boxed{
a^2-2ab+b^2
=
(a-b)^2
}
\]

are perfect-square trinomials.

%%%%%%%%%%%%%%%%%%%%%%%%%%%%%%%%%%%%%%%%%%%%%%%%%%%%%%%%%%%%
\subsection{Completing the Square}
\label{subsec:completing-square-identity}
%%%%%%%%%%%%%%%%%%%%%%%%%%%%%%%%%%%%%%%%%%%%%%%%%%%%%%%%%%%%

For

\[
x^2+bx,
\]

add and subtract

\[
\left(
\frac{b}{2}
\right)^2.
\]

Then

\[
\boxed{
x^2+bx
=
\left(
x+\frac{b}{2}
\right)^2
-
\frac{b^2}{4}.
}
\]

This identity is central to quadratic equations and Gaussian integrals.

%%%%%%%%%%%%%%%%%%%%%%%%%%%%%%%%%%%%%%%%%%%%%%%%%%%%%%%%%%%%
\subsection{Quadratic Completing-the-Square Form}
\label{subsec:quadratic-completing-square}
%%%%%%%%%%%%%%%%%%%%%%%%%%%%%%%%%%%%%%%%%%%%%%%%%%%%%%%%%%%%

For \(a\neq0\),

\[
\boxed{
ax^2+bx+c
=
a
\left(
x+\frac{b}{2a}
\right)^2
+
c
-
\frac{b^2}{4a}.
}
\]

Equivalently,

\[
\boxed{
ax^2+bx+c
=
a
\left(
x+\frac{b}{2a}
\right)^2
-
\frac{b^2-4ac}{4a}.
}
\]

%%%%%%%%%%%%%%%%%%%%%%%%%%%%%%%%%%%%%%%%%%%%%%%%%%%%%%%%%%%%
\subsection{Quadratic Formula}
\label{subsec:quadratic-formula-reference}
%%%%%%%%%%%%%%%%%%%%%%%%%%%%%%%%%%%%%%%%%%%%%%%%%%%%%%%%%%%%

For

\[
ax^2+bx+c=0,
\qquad
a\neq0,
\]

the roots are

\[
\boxed{
x
=
\frac{
-b\pm\sqrt{b^2-4ac}
}{
2a
}.
}
\]

The discriminant is

\[
\boxed{
\Delta
=
b^2-4ac.
}
\]

%%%%%%%%%%%%%%%%%%%%%%%%%%%%%%%%%%%%%%%%%%%%%%%%%%%%%%%%%%%%
\subsection{Discriminant Classification}
\label{subsec:discriminant-classification}
%%%%%%%%%%%%%%%%%%%%%%%%%%%%%%%%%%%%%%%%%%%%%%%%%%%%%%%%%%%%

For real coefficients:

\[
\boxed{
\Delta>0
}
\]

gives two distinct real roots,

\[
\boxed{
\Delta=0
}
\]

gives one repeated real root,

and

\[
\boxed{
\Delta<0
}
\]

gives two nonreal complex conjugate roots.

%%%%%%%%%%%%%%%%%%%%%%%%%%%%%%%%%%%%%%%%%%%%%%%%%%%%%%%%%%%%
\subsection{Vieta's Formulas for Quadratics}
\label{subsec:vieta-quadratic}
%%%%%%%%%%%%%%%%%%%%%%%%%%%%%%%%%%%%%%%%%%%%%%%%%%%%%%%%%%%%

If

\[
ax^2+bx+c
=
a(x-r_1)(x-r_2),
\]

then

\[
\boxed{
r_1+r_2
=
-\frac{b}{a},
}
\]

and

\[
\boxed{
r_1r_2
=
\frac{c}{a}.
}
\]

%%%%%%%%%%%%%%%%%%%%%%%%%%%%%%%%%%%%%%%%%%%%%%%%%%%%%%%%%%%%
\subsection{Polynomial Remainder Identity}
\label{subsec:polynomial-remainder-identity}
%%%%%%%%%%%%%%%%%%%%%%%%%%%%%%%%%%%%%%%%%%%%%%%%%%%%%%%%%%%%

For polynomial \(f(x)\),

\[
\boxed{
f(x)
=
(x-a)q(x)
+
f(a).
}
\]

Thus the remainder when dividing by

\[
x-a
\]

is

\[
\boxed{
f(a).
}
\]

This is the Remainder Theorem.

%%%%%%%%%%%%%%%%%%%%%%%%%%%%%%%%%%%%%%%%%%%%%%%%%%%%%%%%%%%%
\subsection{Factor Theorem}
\label{subsec:factor-theorem-reference}
%%%%%%%%%%%%%%%%%%%%%%%%%%%%%%%%%%%%%%%%%%%%%%%%%%%%%%%%%%%%

The Factor Theorem states

\[
\boxed{
f(a)=0
\iff
(x-a)\text{ is a factor of }f(x).
}
\]

This connects polynomial roots with factorization.

%%%%%%%%%%%%%%%%%%%%%%%%%%%%%%%%%%%%%%%%%%%%%%%%%%%%%%%%%%%%
\subsection{Binomial Theorem}
\label{subsec:binomial-theorem-reference}
%%%%%%%%%%%%%%%%%%%%%%%%%%%%%%%%%%%%%%%%%%%%%%%%%%%%%%%%%%%%

For nonnegative integer \(n\),

\[
\boxed{
(a+b)^n
=
\sum_{k=0}^{n}
\binom{n}{k}
a^{n-k}b^k.
}
\]

The coefficients are binomial coefficients

\[
\boxed{
\binom{n}{k}
=
\frac{n!}{k!(n-k)!}.
}
\]

%%%%%%%%%%%%%%%%%%%%%%%%%%%%%%%%%%%%%%%%%%%%%%%%%%%%%%%%%%%%
\subsection{Binomial Expansion Examples}
\label{subsec:binomial-expansion-examples}
%%%%%%%%%%%%%%%%%%%%%%%%%%%%%%%%%%%%%%%%%%%%%%%%%%%%%%%%%%%%

For \(n=2\),

\[
\boxed{
(a+b)^2
=
a^2+2ab+b^2.
}
\]

For \(n=3\),

\[
\boxed{
(a+b)^3
=
a^3+3a^2b+3ab^2+b^3.
}
\]

For \(n=4\),

\[
\boxed{
(a+b)^4
=
a^4+4a^3b+6a^2b^2+4ab^3+b^4.
}
\]

%%%%%%%%%%%%%%%%%%%%%%%%%%%%%%%%%%%%%%%%%%%%%%%%%%%%%%%%%%%%
\subsection{Pascal's Identity}
\label{subsec:pascals-identity-reference}
%%%%%%%%%%%%%%%%%%%%%%%%%%%%%%%%%%%%%%%%%%%%%%%%%%%%%%%%%%%%

Binomial coefficients satisfy

\[
\boxed{
\binom{n}{k}
=
\binom{n-1}{k-1}
+
\binom{n-1}{k}.
}
\]

This is Pascal's identity.

It generates Pascal's triangle.

%%%%%%%%%%%%%%%%%%%%%%%%%%%%%%%%%%%%%%%%%%%%%%%%%%%%%%%%%%%%
\subsection{Symmetry of Binomial Coefficients}
\label{subsec:binomial-symmetry}
%%%%%%%%%%%%%%%%%%%%%%%%%%%%%%%%%%%%%%%%%%%%%%%%%%%%%%%%%%%%

\[
\boxed{
\binom{n}{k}
=
\binom{n}{n-k}.
}
\]

This reflects the equivalence between choosing \(k\) selected objects and
choosing the \(n-k\) objects left out.

%%%%%%%%%%%%%%%%%%%%%%%%%%%%%%%%%%%%%%%%%%%%%%%%%%%%%%%%%%%%
\subsection{Sum of Binomial Coefficients}
\label{subsec:sum-binomial-coefficients}
%%%%%%%%%%%%%%%%%%%%%%%%%%%%%%%%%%%%%%%%%%%%%%%%%%%%%%%%%%%%

Set

\[
a=b=1
\]

in the binomial theorem:

\[
2^n
=
\sum_{k=0}^{n}
\binom{n}{k}.
\]

Thus

\[
\boxed{
\sum_{k=0}^{n}
\binom{n}{k}
=
2^n.
}
\]

%%%%%%%%%%%%%%%%%%%%%%%%%%%%%%%%%%%%%%%%%%%%%%%%%%%%%%%%%%%%
\subsection{Alternating Binomial Sum}
\label{subsec:alternating-binomial-sum}
%%%%%%%%%%%%%%%%%%%%%%%%%%%%%%%%%%%%%%%%%%%%%%%%%%%%%%%%%%%%

Set

\[
a=1,
\qquad
b=-1.
\]

Then for \(n\geq1\),

\[
0
=
\sum_{k=0}^{n}
(-1)^k
\binom{n}{k}.
\]

Hence

\[
\boxed{
\sum_{k=0}^{n}
(-1)^k
\binom{n}{k}
=
0.
}
\]

%%%%%%%%%%%%%%%%%%%%%%%%%%%%%%%%%%%%%%%%%%%%%%%%%%%%%%%%%%%%
\subsection{Multinomial Theorem}
\label{subsec:multinomial-theorem-reference}
%%%%%%%%%%%%%%%%%%%%%%%%%%%%%%%%%%%%%%%%%%%%%%%%%%%%%%%%%%%%

For nonnegative integer \(n\),

\[
\boxed{
(x_1+\cdots+x_m)^n
=
\sum_{
k_1+\cdots+k_m=n
}
\frac{n!}{
k_1!\cdots k_m!
}
x_1^{k_1}
\cdots
x_m^{k_m}.
}
\]

The coefficients are multinomial coefficients.

%%%%%%%%%%%%%%%%%%%%%%%%%%%%%%%%%%%%%%%%%%%%%%%%%%%%%%%%%%%%
\subsection{Rational Expression Multiplication}
\label{subsec:rational-expression-multiplication}
%%%%%%%%%%%%%%%%%%%%%%%%%%%%%%%%%%%%%%%%%%%%%%%%%%%%%%%%%%%%

For nonzero denominators,

\[
\boxed{
\frac{a}{b}
\frac{c}{d}
=
\frac{ac}{bd}.
}
\]

%%%%%%%%%%%%%%%%%%%%%%%%%%%%%%%%%%%%%%%%%%%%%%%%%%%%%%%%%%%%
\subsection{Rational Expression Division}
\label{subsec:rational-expression-division}
%%%%%%%%%%%%%%%%%%%%%%%%%%%%%%%%%%%%%%%%%%%%%%%%%%%%%%%%%%%%

For nonzero \(b,c,d\),

\[
\boxed{
\frac{a/b}{c/d}
=
\frac{a}{b}
\frac{d}{c}
=
\frac{ad}{bc}.
}
\]

Division by a fraction means multiplication by its reciprocal.

%%%%%%%%%%%%%%%%%%%%%%%%%%%%%%%%%%%%%%%%%%%%%%%%%%%%%%%%%%%%
\subsection{Adding Rational Expressions}
\label{subsec:rational-expression-addition}
%%%%%%%%%%%%%%%%%%%%%%%%%%%%%%%%%%%%%%%%%%%%%%%%%%%%%%%%%%%%

For nonzero denominators,

\[
\boxed{
\frac{a}{b}
+
\frac{c}{d}
=
\frac{ad+bc}{bd}.
}
\]

Similarly,

\[
\boxed{
\frac{a}{b}
-
\frac{c}{d}
=
\frac{ad-bc}{bd}.
}
\]

%%%%%%%%%%%%%%%%%%%%%%%%%%%%%%%%%%%%%%%%%%%%%%%%%%%%%%%%%%%%
\subsection{Cancellation Rule}
\label{subsec:cancellation-rule}
%%%%%%%%%%%%%%%%%%%%%%%%%%%%%%%%%%%%%%%%%%%%%%%%%%%%%%%%%%%%

If \(c\neq0\),

\[
\boxed{
\frac{ac}{bc}
=
\frac{a}{b}.
}
\]

Only common factors may be canceled.

Terms connected by addition cannot generally be canceled.

%%%%%%%%%%%%%%%%%%%%%%%%%%%%%%%%%%%%%%%%%%%%%%%%%%%%%%%%%%%%
\subsection{Common False Cancellation}
\label{subsec:false-cancellation}
%%%%%%%%%%%%%%%%%%%%%%%%%%%%%%%%%%%%%%%%%%%%%%%%%%%%%%%%%%%%

In general,

\[
\boxed{
\frac{a+b}{a}
\neq
b.
}
\]

Instead,

\[
\boxed{
\frac{a+b}{a}
=
1+\frac{b}{a},
\qquad
a\neq0.
}
\]

Cancellation applies to factors, not terms.

%%%%%%%%%%%%%%%%%%%%%%%%%%%%%%%%%%%%%%%%%%%%%%%%%%%%%%%%%%%%
\subsection{Complex Fractions}
\label{subsec:complex-fractions-reference}
%%%%%%%%%%%%%%%%%%%%%%%%%%%%%%%%%%%%%%%%%%%%%%%%%%%%%%%%%%%%

A complex fraction contains fractions in its numerator or denominator.

For example,

\[
\frac{
\frac1a+\frac1b
}{
\frac1c
}.
\]

Multiply numerator and denominator by an appropriate common denominator to
simplify.

%%%%%%%%%%%%%%%%%%%%%%%%%%%%%%%%%%%%%%%%%%%%%%%%%%%%%%%%%%%%
\subsection{Worked Example: Complex Fraction}
\label{subsec:worked-complex-fraction}
%%%%%%%%%%%%%%%%%%%%%%%%%%%%%%%%%%%%%%%%%%%%%%%%%%%%%%%%%%%%

\begin{workedexamplebox}{Simplify a Nested Fraction}

Simplify

\[
\frac{
\frac1x+\frac1y
}{
\frac1{xy}
}.
\]

First,

\[
\frac1x+\frac1y
=
\frac{x+y}{xy}.
\]

Then

\[
\frac{
\frac{x+y}{xy}
}{
\frac1{xy}
}
=
x+y,
\]

assuming

\[
x\neq0,
\qquad
y\neq0.
\]

\begin{answerbox}
\[
\boxed{
\frac{
\frac1x+\frac1y
}{
\frac1{xy}
}
=
x+y.
}
\]
\end{answerbox}

\end{workedexamplebox}

%%%%%%%%%%%%%%%%%%%%%%%%%%%%%%%%%%%%%%%%%%%%%%%%%%%%%%%%%%%%
\subsection{Reciprocal of a Product}
\label{subsec:reciprocal-product}
%%%%%%%%%%%%%%%%%%%%%%%%%%%%%%%%%%%%%%%%%%%%%%%%%%%%%%%%%%%%

For nonzero \(a,b\),

\[
\boxed{
\frac{1}{ab}
=
\frac1a\frac1b.
}
\]

%%%%%%%%%%%%%%%%%%%%%%%%%%%%%%%%%%%%%%%%%%%%%%%%%%%%%%%%%%%%
\subsection{Reciprocal of a Quotient}
\label{subsec:reciprocal-quotient}
%%%%%%%%%%%%%%%%%%%%%%%%%%%%%%%%%%%%%%%%%%%%%%%%%%%%%%%%%%%%

For nonzero \(a,b\),

\[
\boxed{
\frac{1}{a/b}
=
\frac{b}{a}.
}
\]

%%%%%%%%%%%%%%%%%%%%%%%%%%%%%%%%%%%%%%%%%%%%%%%%%%%%%%%%%%%%
\subsection{Absolute Value Product Rule}
\label{subsec:absolute-value-product}
%%%%%%%%%%%%%%%%%%%%%%%%%%%%%%%%%%%%%%%%%%%%%%%%%%%%%%%%%%%%

For real or complex numbers,

\[
\boxed{
|ab|
=
|a||b|.
}
\]

%%%%%%%%%%%%%%%%%%%%%%%%%%%%%%%%%%%%%%%%%%%%%%%%%%%%%%%%%%%%
\subsection{Absolute Value Quotient Rule}
\label{subsec:absolute-value-quotient}
%%%%%%%%%%%%%%%%%%%%%%%%%%%%%%%%%%%%%%%%%%%%%%%%%%%%%%%%%%%%

If \(b\neq0\),

\[
\boxed{
\left|
\frac{a}{b}
\right|
=
\frac{|a|}{|b|}.
}
\]

%%%%%%%%%%%%%%%%%%%%%%%%%%%%%%%%%%%%%%%%%%%%%%%%%%%%%%%%%%%%
\subsection{Absolute Value and Powers}
\label{subsec:absolute-value-powers}
%%%%%%%%%%%%%%%%%%%%%%%%%%%%%%%%%%%%%%%%%%%%%%%%%%%%%%%%%%%%

For integer \(n\geq0\),

\[
\boxed{
|a^n|
=
|a|^n.
}
\]

In particular,

\[
\boxed{
a^2
=
|a|^2
}
\]

for real \(a\).

%%%%%%%%%%%%%%%%%%%%%%%%%%%%%%%%%%%%%%%%%%%%%%%%%%%%%%%%%%%%
\subsection{Triangle Inequality}
\label{subsec:triangle-inequality-algebraic}
%%%%%%%%%%%%%%%%%%%%%%%%%%%%%%%%%%%%%%%%%%%%%%%%%%%%%%%%%%%%

For real or complex \(a,b\),

\[
\boxed{
|a+b|
\leq
|a|+|b|.
}
\]

This is the triangle inequality.

%%%%%%%%%%%%%%%%%%%%%%%%%%%%%%%%%%%%%%%%%%%%%%%%%%%%%%%%%%%%
\subsection{Reverse Triangle Inequality}
\label{subsec:reverse-triangle-inequality}
%%%%%%%%%%%%%%%%%%%%%%%%%%%%%%%%%%%%%%%%%%%%%%%%%%%%%%%%%%%%

A related identity is

\[
\boxed{
\bigl||a|-|b|\bigr|
\leq
|a-b|.
}
\]

This is useful throughout analysis.

%%%%%%%%%%%%%%%%%%%%%%%%%%%%%%%%%%%%%%%%%%%%%%%%%%%%%%%%%%%%
\subsection{Common False Absolute-Value Identity}
\label{subsec:false-absolute-value-identity}
%%%%%%%%%%%%%%%%%%%%%%%%%%%%%%%%%%%%%%%%%%%%%%%%%%%%%%%%%%%%

In general,

\[
\boxed{
|a+b|
\neq
|a|+|b|.
}
\]

Equality requires additional conditions.

For real numbers, one sufficient condition is that \(a\) and \(b\) have the
same sign.

%%%%%%%%%%%%%%%%%%%%%%%%%%%%%%%%%%%%%%%%%%%%%%%%%%%%%%%%%%%%
\subsection{Complex Conjugate Identities}
\label{subsec:complex-conjugate-identities}
%%%%%%%%%%%%%%%%%%%%%%%%%%%%%%%%%%%%%%%%%%%%%%%%%%%%%%%%%%%%

For complex numbers \(z,w\),

\[
\boxed{
\overline{z+w}
=
\overline z+\overline w,
}
\]

\[
\boxed{
\overline{zw}
=
\overline z\,\overline w,
}
\]

and

\[
\boxed{
\overline{\overline z}
=
z.
}
\]

%%%%%%%%%%%%%%%%%%%%%%%%%%%%%%%%%%%%%%%%%%%%%%%%%%%%%%%%%%%%
\subsection{Complex Modulus Identity}
\label{subsec:complex-modulus-identity}
%%%%%%%%%%%%%%%%%%%%%%%%%%%%%%%%%%%%%%%%%%%%%%%%%%%%%%%%%%%%

For complex \(z\),

\[
\boxed{
|z|^2
=
z\overline z.
}
\]

If

\[
z=a+bi,
\]

then

\[
z\overline z
=
(a+bi)(a-bi)
=
a^2+b^2.
\]

%%%%%%%%%%%%%%%%%%%%%%%%%%%%%%%%%%%%%%%%%%%%%%%%%%%%%%%%%%%%
\subsection{Real and Imaginary Parts From Conjugates}
\label{subsec:real-imag-conjugate-identities}
%%%%%%%%%%%%%%%%%%%%%%%%%%%%%%%%%%%%%%%%%%%%%%%%%%%%%%%%%%%%

For complex \(z\),

\[
\boxed{
\operatorname{Re}(z)
=
\frac{
z+\overline z
}{2},
}
\]

and

\[
\boxed{
\operatorname{Im}(z)
=
\frac{
z-\overline z
}{
2i
}.
}
\]

%%%%%%%%%%%%%%%%%%%%%%%%%%%%%%%%%%%%%%%%%%%%%%%%%%%%%%%%%%%%
\subsection{Euler Formula}
\label{subsec:euler-formula-algebraic-identities}
%%%%%%%%%%%%%%%%%%%%%%%%%%%%%%%%%%%%%%%%%%%%%%%%%%%%%%%%%%%%

For real \(\theta\),

\[
\boxed{
e^{i\theta}
=
\cos\theta+i\sin\theta.
}
\]

Therefore,

\[
\boxed{
e^{-i\theta}
=
\cos\theta-i\sin\theta.
}
\]

%%%%%%%%%%%%%%%%%%%%%%%%%%%%%%%%%%%%%%%%%%%%%%%%%%%%%%%%%%%%
\subsection{Sine and Cosine From Exponentials}
\label{subsec:sine-cosine-exponential-identities}
%%%%%%%%%%%%%%%%%%%%%%%%%%%%%%%%%%%%%%%%%%%%%%%%%%%%%%%%%%%%

Adding Euler formulas gives

\[
\boxed{
\cos\theta
=
\frac{
e^{i\theta}+e^{-i\theta}
}{2}.
}
\]

Subtracting gives

\[
\boxed{
\sin\theta
=
\frac{
e^{i\theta}-e^{-i\theta}
}{
2i
}.
}
\]

%%%%%%%%%%%%%%%%%%%%%%%%%%%%%%%%%%%%%%%%%%%%%%%%%%%%%%%%%%%%
\subsection{Hyperbolic Identities}
\label{subsec:hyperbolic-identities-reference}
%%%%%%%%%%%%%%%%%%%%%%%%%%%%%%%%%%%%%%%%%%%%%%%%%%%%%%%%%%%%

The hyperbolic functions satisfy

\[
\boxed{
\cosh x
=
\frac{
e^x+e^{-x}
}{2},
}
\]

\[
\boxed{
\sinh x
=
\frac{
e^x-e^{-x}
}{2}.
}
\]

Therefore,

\[
\boxed{
\cosh^2x-\sinh^2x=1.
}
\]

%%%%%%%%%%%%%%%%%%%%%%%%%%%%%%%%%%%%%%%%%%%%%%%%%%%%%%%%%%%%
\subsection{Difference of Reciprocal Squares}
\label{subsec:reciprocal-difference-identity}
%%%%%%%%%%%%%%%%%%%%%%%%%%%%%%%%%%%%%%%%%%%%%%%%%%%%%%%%%%%%

For nonzero \(a,b\),

\[
\boxed{
\frac1{a^2}
-
\frac1{b^2}
=
\frac{
b^2-a^2
}{
a^2b^2
}.
}
\]

Factoring the numerator,

\[
\boxed{
\frac1{a^2}
-
\frac1{b^2}
=
\frac{
(b-a)(b+a)
}{
a^2b^2
}.
}
\]

%%%%%%%%%%%%%%%%%%%%%%%%%%%%%%%%%%%%%%%%%%%%%%%%%%%%%%%%%%%%
\subsection{Difference of Reciprocals}
\label{subsec:difference-reciprocals}
%%%%%%%%%%%%%%%%%%%%%%%%%%%%%%%%%%%%%%%%%%%%%%%%%%%%%%%%%%%%

For nonzero \(a,b\),

\[
\boxed{
\frac1a-\frac1b
=
\frac{b-a}{ab}.
}
\]

Similarly,

\[
\boxed{
\frac1a+\frac1b
=
\frac{a+b}{ab}.
}
\]

%%%%%%%%%%%%%%%%%%%%%%%%%%%%%%%%%%%%%%%%%%%%%%%%%%%%%%%%%%%%
\subsection{Useful Fraction Identity}
\label{subsec:useful-fraction-identity}
%%%%%%%%%%%%%%%%%%%%%%%%%%%%%%%%%%%%%%%%%%%%%%%%%%%%%%%%%%%%

For nonzero \(x\),

\[
\boxed{
x+\frac1x
}
\]

has square

\[
\boxed{
\left(
x+\frac1x
\right)^2
=
x^2+2+\frac1{x^2}.
}
\]

Therefore,

\[
\boxed{
x^2+\frac1{x^2}
=
\left(
x+\frac1x
\right)^2-2.
}
\]

%%%%%%%%%%%%%%%%%%%%%%%%%%%%%%%%%%%%%%%%%%%%%%%%%%%%%%%%%%%%
\subsection{Related Reciprocal Identity}
\label{subsec:related-reciprocal-identity}
%%%%%%%%%%%%%%%%%%%%%%%%%%%%%%%%%%%%%%%%%%%%%%%%%%%%%%%%%%%%

Likewise,

\[
\boxed{
\left(
x-\frac1x
\right)^2
=
x^2-2+\frac1{x^2}.
}
\]

Hence,

\[
\boxed{
x^2+\frac1{x^2}
=
\left(
x-\frac1x
\right)^2+2.
}
\]

%%%%%%%%%%%%%%%%%%%%%%%%%%%%%%%%%%%%%%%%%%%%%%%%%%%%%%%%%%%%
\subsection{Cubic Reciprocal Identity}
\label{subsec:cubic-reciprocal-identity}
%%%%%%%%%%%%%%%%%%%%%%%%%%%%%%%%%%%%%%%%%%%%%%%%%%%%%%%%%%%%

Expanding,

\[
\left(
x+\frac1x
\right)^3
=
x^3
+
3x
+
\frac3x
+
\frac1{x^3}.
\]

Therefore,

\[
\boxed{
x^3+\frac1{x^3}
=
\left(
x+\frac1x
\right)^3
-
3
\left(
x+\frac1x
\right).
}
\]

%%%%%%%%%%%%%%%%%%%%%%%%%%%%%%%%%%%%%%%%%%%%%%%%%%%%%%%%%%%%
\subsection{Arithmetic Sequence Sum}
\label{subsec:arithmetic-sequence-sum-identity}
%%%%%%%%%%%%%%%%%%%%%%%%%%%%%%%%%%%%%%%%%%%%%%%%%%%%%%%%%%%%

For arithmetic sequence

\[
a,
\quad
a+d,
\quad
\ldots,
\quad
a+(n-1)d,
\]

the sum is

\[
\boxed{
S_n
=
\frac{n}{2}
\left[
2a+(n-1)d
\right].
}
\]

Equivalently,

\[
\boxed{
S_n
=
\frac{n}{2}
(a_1+a_n).
}
\]

%%%%%%%%%%%%%%%%%%%%%%%%%%%%%%%%%%%%%%%%%%%%%%%%%%%%%%%%%%%%
\subsection{Geometric Sequence Sum}
\label{subsec:geometric-sequence-sum-identity}
%%%%%%%%%%%%%%%%%%%%%%%%%%%%%%%%%%%%%%%%%%%%%%%%%%%%%%%%%%%%

For

\[
r\neq1,
\]

\[
\boxed{
1+r+r^2+\cdots+r^{n-1}
=
\frac{
1-r^n
}{
1-r
}.
}
\]

Equivalently,

\[
\boxed{
a+ar+\cdots+ar^{n-1}
=
a
\frac{
1-r^n
}{
1-r
}.
}
\]

%%%%%%%%%%%%%%%%%%%%%%%%%%%%%%%%%%%%%%%%%%%%%%%%%%%%%%%%%%%%
\subsection{Infinite Geometric Series}
\label{subsec:infinite-geometric-series-reference}
%%%%%%%%%%%%%%%%%%%%%%%%%%%%%%%%%%%%%%%%%%%%%%%%%%%%%%%%%%%%

If

\[
|r|<1,
\]

then

\[
\boxed{
1+r+r^2+r^3+\cdots
=
\frac{1}{1-r}.
}
\]

More generally,

\[
\boxed{
\sum_{n=0}^{\infty}
ar^n
=
\frac{a}{1-r}.
}
\]

%%%%%%%%%%%%%%%%%%%%%%%%%%%%%%%%%%%%%%%%%%%%%%%%%%%%%%%%%%%%
\subsection{Sum of First Integers}
\label{subsec:sum-first-integers}
%%%%%%%%%%%%%%%%%%%%%%%%%%%%%%%%%%%%%%%%%%%%%%%%%%%%%%%%%%%%

\[
\boxed{
1+2+\cdots+n
=
\frac{n(n+1)}{2}.
}
\]

Equivalently,

\[
\boxed{
\sum_{k=1}^{n}k
=
\frac{n(n+1)}{2}.
}
\]

%%%%%%%%%%%%%%%%%%%%%%%%%%%%%%%%%%%%%%%%%%%%%%%%%%%%%%%%%%%%
\subsection{Sum of First Squares}
\label{subsec:sum-first-squares}
%%%%%%%%%%%%%%%%%%%%%%%%%%%%%%%%%%%%%%%%%%%%%%%%%%%%%%%%%%%%

\[
\boxed{
\sum_{k=1}^{n}k^2
=
\frac{
n(n+1)(2n+1)
}{6}.
}
\]

%%%%%%%%%%%%%%%%%%%%%%%%%%%%%%%%%%%%%%%%%%%%%%%%%%%%%%%%%%%%
\subsection{Sum of First Cubes}
\label{subsec:sum-first-cubes}
%%%%%%%%%%%%%%%%%%%%%%%%%%%%%%%%%%%%%%%%%%%%%%%%%%%%%%%%%%%%

\[
\boxed{
\sum_{k=1}^{n}k^3
=
\left[
\frac{n(n+1)}{2}
\right]^2.
}
\]

Thus

\[
\boxed{
1^3+2^3+\cdots+n^3
=
(1+2+\cdots+n)^2.
}
\]

%%%%%%%%%%%%%%%%%%%%%%%%%%%%%%%%%%%%%%%%%%%%%%%%%%%%%%%%%%%%
\subsection{Telescoping Identity}
\label{subsec:telescoping-identity-reference}
%%%%%%%%%%%%%%%%%%%%%%%%%%%%%%%%%%%%%%%%%%%%%%%%%%%%%%%%%%%%

A useful partial-fraction identity is

\[
\boxed{
\frac{1}{k(k+1)}
=
\frac1k
-
\frac1{k+1}.
}
\]

Therefore,

\[
\sum_{k=1}^{n}
\frac{1}{k(k+1)}
\]

telescopes:

\[
\boxed{
\sum_{k=1}^{n}
\frac{1}{k(k+1)}
=
1-\frac{1}{n+1}.
}
\]

%%%%%%%%%%%%%%%%%%%%%%%%%%%%%%%%%%%%%%%%%%%%%%%%%%%%%%%%%%%%
\subsection{Partial Fraction Identity}
\label{subsec:basic-partial-fraction-identity}
%%%%%%%%%%%%%%%%%%%%%%%%%%%%%%%%%%%%%%%%%%%%%%%%%%%%%%%%%%%%

For distinct \(a,b\),

\[
\boxed{
\frac{1}{
(x-a)(x-b)
}
=
\frac{1}{a-b}
\left(
\frac{1}{x-a}
-
\frac{1}{x-b}
\right).
}
\]

This identity is useful in integration and series.

%%%%%%%%%%%%%%%%%%%%%%%%%%%%%%%%%%%%%%%%%%%%%%%%%%%%%%%%%%%%
\subsection{Difference Quotient Identity}
\label{subsec:difference-quotient-polynomial}
%%%%%%%%%%%%%%%%%%%%%%%%%%%%%%%%%%%%%%%%%%%%%%%%%%%%%%%%%%%%

For positive integer \(n\),

\[
\boxed{
\frac{
x^n-a^n
}{
x-a
}
=
x^{n-1}
+
x^{n-2}a
+
\cdots
+
xa^{n-2}
+
a^{n-1},
}
\]

provided

\[
x\neq a.
\]

This is directly related to the factorization of \(x^n-a^n\).

%%%%%%%%%%%%%%%%%%%%%%%%%%%%%%%%%%%%%%%%%%%%%%%%%%%%%%%%%%%%
\subsection{Worked Example: Difference Quotient}
\label{subsec:worked-difference-quotient}
%%%%%%%%%%%%%%%%%%%%%%%%%%%%%%%%%%%%%%%%%%%%%%%%%%%%%%%%%%%%

\begin{workedexamplebox}{Simplify a Polynomial Difference Quotient}

Simplify

\[
\frac{x^3-a^3}{x-a}.
\]

Factor the numerator:

\[
x^3-a^3
=
(x-a)(x^2+ax+a^2).
\]

Therefore, when \(x\neq a\),

\begin{answerbox}
\[
\boxed{
\frac{x^3-a^3}{x-a}
=
x^2+ax+a^2.
}
\]
\end{answerbox}

\end{workedexamplebox}

%%%%%%%%%%%%%%%%%%%%%%%%%%%%%%%%%%%%%%%%%%%%%%%%%%%%%%%%%%%%
\subsection{Difference Quotient for Squares}
\label{subsec:difference-quotient-squares}
%%%%%%%%%%%%%%%%%%%%%%%%%%%%%%%%%%%%%%%%%%%%%%%%%%%%%%%%%%%%

A particularly important identity is

\[
\boxed{
\frac{
(x+h)^2-x^2
}{
h
}
=
2x+h,
\qquad
h\neq0.
}
\]

This appears directly in derivative calculations.

%%%%%%%%%%%%%%%%%%%%%%%%%%%%%%%%%%%%%%%%%%%%%%%%%%%%%%%%%%%%
\subsection{Difference Quotient for Powers}
\label{subsec:difference-quotient-powers}
%%%%%%%%%%%%%%%%%%%%%%%%%%%%%%%%%%%%%%%%%%%%%%%%%%%%%%%%%%%%

By the binomial theorem,

\[
(x+h)^n
=
x^n
+
nx^{n-1}h
+
\binom n2x^{n-2}h^2
+
\cdots
+
h^n.
\]

Therefore,

\[
\boxed{
\frac{
(x+h)^n-x^n
}{
h
}
=
nx^{n-1}
+
\binom n2x^{n-2}h
+
\cdots
+
h^{n-1}.
}
\]

As \(h\to0\),

\[
\boxed{
\frac{d}{dx}x^n
=
nx^{n-1}.
}
\]

%%%%%%%%%%%%%%%%%%%%%%%%%%%%%%%%%%%%%%%%%%%%%%%%%%%%%%%%%%%%
\subsection{Difference of Function Values}
\label{subsec:function-difference-factorization}
%%%%%%%%%%%%%%%%%%%%%%%%%%%%%%%%%%%%%%%%%%%%%%%%%%%%%%%%%%%%

For quadratic

\[
f(x)=ax^2+bx+c,
\]

\[
f(x)-f(y)
=
a(x^2-y^2)+b(x-y).
\]

Thus

\[
\boxed{
f(x)-f(y)
=
(x-y)
\left[
a(x+y)+b
\right].
}
\]

This type of factorization is useful in injectivity and continuity arguments.

%%%%%%%%%%%%%%%%%%%%%%%%%%%%%%%%%%%%%%%%%%%%%%%%%%%%%%%%%%%%
\subsection{Completing Squares in Two Variables}
\label{subsec:completing-squares-two-variables}
%%%%%%%%%%%%%%%%%%%%%%%%%%%%%%%%%%%%%%%%%%%%%%%%%%%%%%%%%%%%

For real \(x,y\),

\[
\boxed{
x^2+y^2
=
(x-y)^2+2xy
}
\]

and

\[
\boxed{
x^2+y^2
=
(x+y)^2-2xy.
}
\]

These identities are often used in inequalities.

%%%%%%%%%%%%%%%%%%%%%%%%%%%%%%%%%%%%%%%%%%%%%%%%%%%%%%%%%%%%
\subsection{Useful Quadratic Identity}
\label{subsec:useful-quadratic-identity}
%%%%%%%%%%%%%%%%%%%%%%%%%%%%%%%%%%%%%%%%%%%%%%%%%%%%%%%%%%%%

\[
\boxed{
(a+b)^2+(a-b)^2
=
2a^2+2b^2.
}
\]

Similarly,

\[
\boxed{
(a+b)^2-(a-b)^2
=
4ab.
}
\]

%%%%%%%%%%%%%%%%%%%%%%%%%%%%%%%%%%%%%%%%%%%%%%%%%%%%%%%%%%%%
\subsection{Parallelogram Identity}
\label{subsec:parallelogram-identity-reference}
%%%%%%%%%%%%%%%%%%%%%%%%%%%%%%%%%%%%%%%%%%%%%%%%%%%%%%%%%%%%

In Euclidean or inner-product spaces,

\[
\boxed{
\|x+y\|^2
+
\|x-y\|^2
=
2\|x\|^2
+
2\|y\|^2.
}
\]

This is the parallelogram identity.

It characterizes norms induced by inner products.

%%%%%%%%%%%%%%%%%%%%%%%%%%%%%%%%%%%%%%%%%%%%%%%%%%%%%%%%%%%%
\subsection{Dot Product Expansion}
\label{subsec:dot-product-expansion}
%%%%%%%%%%%%%%%%%%%%%%%%%%%%%%%%%%%%%%%%%%%%%%%%%%%%%%%%%%%%

For vectors \(u,v\),

\[
\boxed{
\|u+v\|^2
=
\|u\|^2
+
2u\cdot v
+
\|v\|^2.
}
\]

Similarly,

\[
\boxed{
\|u-v\|^2
=
\|u\|^2
-
2u\cdot v
+
\|v\|^2.
}
\]

%%%%%%%%%%%%%%%%%%%%%%%%%%%%%%%%%%%%%%%%%%%%%%%%%%%%%%%%%%%%
\subsection{Polarization Identity}
\label{subsec:polarization-identity-reference}
%%%%%%%%%%%%%%%%%%%%%%%%%%%%%%%%%%%%%%%%%%%%%%%%%%%%%%%%%%%%

For a real inner-product space,

\[
\boxed{
\langle x,y\rangle
=
\frac14
\left(
\|x+y\|^2
-
\|x-y\|^2
\right).
}
\]

This reconstructs the inner product from the norm.

%%%%%%%%%%%%%%%%%%%%%%%%%%%%%%%%%%%%%%%%%%%%%%%%%%%%%%%%%%%%
\subsection{Determinant Identity for \(2\times2\) Matrices}
\label{subsec:determinant-two-by-two-identity}
%%%%%%%%%%%%%%%%%%%%%%%%%%%%%%%%%%%%%%%%%%%%%%%%%%%%%%%%%%%%

For

\[
A
=
\begin{bmatrix}
a&b\\
c&d
\end{bmatrix},
\]

\[
\boxed{
\det(A)
=
ad-bc.
}
\]

If

\[
\det(A)\neq0,
\]

then

\[
\boxed{
A^{-1}
=
\frac{1}{ad-bc}
\begin{bmatrix}
d&-b\\
-c&a
\end{bmatrix}.
}
\]

%%%%%%%%%%%%%%%%%%%%%%%%%%%%%%%%%%%%%%%%%%%%%%%%%%%%%%%%%%%%
\subsection{Trace Identities}
\label{subsec:trace-identities-reference}
%%%%%%%%%%%%%%%%%%%%%%%%%%%%%%%%%%%%%%%%%%%%%%%%%%%%%%%%%%%%

For compatible square matrices,

\[
\boxed{
\operatorname{tr}(A+B)
=
\operatorname{tr}(A)
+
\operatorname{tr}(B).
}
\]

Also,

\[
\boxed{
\operatorname{tr}(AB)
=
\operatorname{tr}(BA).
}
\]

More generally, cyclic permutations satisfy

\[
\boxed{
\operatorname{tr}(ABC)
=
\operatorname{tr}(BCA)
=
\operatorname{tr}(CAB).
}
\]

%%%%%%%%%%%%%%%%%%%%%%%%%%%%%%%%%%%%%%%%%%%%%%%%%%%%%%%%%%%%
\subsection{Transpose Identities}
\label{subsec:transpose-identities-reference}
%%%%%%%%%%%%%%%%%%%%%%%%%%%%%%%%%%%%%%%%%%%%%%%%%%%%%%%%%%%%

For compatible matrices,

\[
\boxed{
(A+B)^T
=
A^T+B^T,
}
\]

\[
\boxed{
(AB)^T
=
B^TA^T,
}
\]

and

\[
\boxed{
(A^T)^T
=
A.
}
\]

%%%%%%%%%%%%%%%%%%%%%%%%%%%%%%%%%%%%%%%%%%%%%%%%%%%%%%%%%%%%
\subsection{Inverse Matrix Identities}
\label{subsec:inverse-matrix-identities}
%%%%%%%%%%%%%%%%%%%%%%%%%%%%%%%%%%%%%%%%%%%%%%%%%%%%%%%%%%%%

For invertible matrices,

\[
\boxed{
(AB)^{-1}
=
B^{-1}A^{-1}.
}
\]

Also,

\[
\boxed{
(A^T)^{-1}
=
(A^{-1})^T.
}
\]

Note the reversed order in the inverse of a product.

%%%%%%%%%%%%%%%%%%%%%%%%%%%%%%%%%%%%%%%%%%%%%%%%%%%%%%%%%%%%
\subsection{Determinant Product Identity}
\label{subsec:determinant-product-identity}
%%%%%%%%%%%%%%%%%%%%%%%%%%%%%%%%%%%%%%%%%%%%%%%%%%%%%%%%%%%%

For square matrices,

\[
\boxed{
\det(AB)
=
\det(A)\det(B).
}
\]

If \(A\) is invertible,

\[
\boxed{
\det(A^{-1})
=
\frac{1}{\det(A)}.
}
\]

%%%%%%%%%%%%%%%%%%%%%%%%%%%%%%%%%%%%%%%%%%%%%%%%%%%%%%%%%%%%
\subsection{Geometric-Series Matrix Identity}
\label{subsec:matrix-geometric-series}
%%%%%%%%%%%%%%%%%%%%%%%%%%%%%%%%%%%%%%%%%%%%%%%%%%%%%%%%%%%%

For a matrix \(A\),

\[
\boxed{
(I-A)
(I+A+A^2+\cdots+A^{n-1})
=
I-A^n.
}
\]

If the relevant inverse exists,

\[
\boxed{
I+A+\cdots+A^{n-1}
=
(I-A^n)(I-A)^{-1}.
}
\]

%%%%%%%%%%%%%%%%%%%%%%%%%%%%%%%%%%%%%%%%%%%%%%%%%%%%%%%%%%%%
\subsection{Infinite Matrix Geometric Series}
\label{subsec:infinite-matrix-geometric-series}
%%%%%%%%%%%%%%%%%%%%%%%%%%%%%%%%%%%%%%%%%%%%%%%%%%%%%%%%%%%%

If the spectral radius satisfies

\[
\rho(A)<1,
\]

then

\[
\boxed{
(I-A)^{-1}
=
I+A+A^2+A^3+\cdots.
}
\]

This is the Neumann series.

%%%%%%%%%%%%%%%%%%%%%%%%%%%%%%%%%%%%%%%%%%%%%%%%%%%%%%%%%%%%
\subsection{Common Domain Restrictions}
\label{subsec:algebra-domain-restrictions}
%%%%%%%%%%%%%%%%%%%%%%%%%%%%%%%%%%%%%%%%%%%%%%%%%%%%%%%%%%%%

Algebraic simplification must preserve the domain.

Examples include:

\[
\boxed{
\frac{x^2-1}{x-1}
=
x+1
}
\]

only when

\[
\boxed{
x\neq1.
}
\]

The simplified expression is defined at \(x=1\), but the original expression
is not.

%%%%%%%%%%%%%%%%%%%%%%%%%%%%%%%%%%%%%%%%%%%%%%%%%%%%%%%%%%%%
\subsection{Equivalent Expressions Versus Equal Functions}
\label{subsec:equivalent-expressions-functions}
%%%%%%%%%%%%%%%%%%%%%%%%%%%%%%%%%%%%%%%%%%%%%%%%%%%%%%%%%%%%

The expressions

\[
\frac{x^2-1}{x-1}
\]

and

\[
x+1
\]

agree wherever the first is defined.

But as functions with their natural domains, they are not identical because
their domains differ.

Thus algebraic simplification should track excluded values.

%%%%%%%%%%%%%%%%%%%%%%%%%%%%%%%%%%%%%%%%%%%%%%%%%%%%%%%%%%%%
\subsection{Squaring Both Sides}
\label{subsec:squaring-both-sides-warning}
%%%%%%%%%%%%%%%%%%%%%%%%%%%%%%%%%%%%%%%%%%%%%%%%%%%%%%%%%%%%

If

\[
a=b,
\]

then

\[
a^2=b^2.
\]

But the reverse implication fails:

\[
a^2=b^2
\]

implies

\[
\boxed{
a=\pm b.
}
\]

Squaring an equation can introduce extraneous solutions.

%%%%%%%%%%%%%%%%%%%%%%%%%%%%%%%%%%%%%%%%%%%%%%%%%%%%%%%%%%%%
\subsection{Taking Square Roots}
\label{subsec:taking-square-roots-warning}
%%%%%%%%%%%%%%%%%%%%%%%%%%%%%%%%%%%%%%%%%%%%%%%%%%%%%%%%%%%%

From

\[
x^2=a,
\qquad
a\geq0,
\]

one obtains

\[
\boxed{
x=\pm\sqrt a.
}
\]

The symbol

\[
\sqrt a
\]

by itself denotes the principal nonnegative square root.

%%%%%%%%%%%%%%%%%%%%%%%%%%%%%%%%%%%%%%%%%%%%%%%%%%%%%%%%%%%%
\subsection{Multiplying by an Expression}
\label{subsec:multiplying-expression-domain-warning}
%%%%%%%%%%%%%%%%%%%%%%%%%%%%%%%%%%%%%%%%%%%%%%%%%%%%%%%%%%%%

Multiplying both sides by an expression that could equal zero requires care.

For example, multiplying by \(x\) can change logical equivalence if \(x=0\)
is allowed.

Always record conditions such as

\[
\boxed{
x\neq0.
}
\]

%%%%%%%%%%%%%%%%%%%%%%%%%%%%%%%%%%%%%%%%%%%%%%%%%%%%%%%%%%%%
\subsection{Dividing by an Expression}
\label{subsec:dividing-expression-warning}
%%%%%%%%%%%%%%%%%%%%%%%%%%%%%%%%%%%%%%%%%%%%%%%%%%%%%%%%%%%%

Division by an expression is valid only when that expression is nonzero.

For example,

\[
ax=bx
\]

does not imply

\[
a=b
\]

unless

\[
\boxed{
x\neq0.
}
\]

If \(x=0\), both sides are equal regardless of \(a\) and \(b\).

%%%%%%%%%%%%%%%%%%%%%%%%%%%%%%%%%%%%%%%%%%%%%%%%%%%%%%%%%%%%
\subsection{Common False Exponent Identity}
\label{subsec:false-exponent-sum}
%%%%%%%%%%%%%%%%%%%%%%%%%%%%%%%%%%%%%%%%%%%%%%%%%%%%%%%%%%%%

In general,

\[
\boxed{
(a+b)^n
\neq
a^n+b^n.
}
\]

For \(n=2\),

\[
(a+b)^2
=
a^2+2ab+b^2.
\]

The missing cross terms are essential.

%%%%%%%%%%%%%%%%%%%%%%%%%%%%%%%%%%%%%%%%%%%%%%%%%%%%%%%%%%%%
\subsection{Common False Radical Identity}
\label{subsec:false-radical-sum}
%%%%%%%%%%%%%%%%%%%%%%%%%%%%%%%%%%%%%%%%%%%%%%%%%%%%%%%%%%%%

In general,

\[
\boxed{
\sqrt{a+b}
\neq
\sqrt a+\sqrt b.
}
\]

For example,

\[
\sqrt{1+1}
=
\sqrt2,
\]

while

\[
\sqrt1+\sqrt1=2.
\]

%%%%%%%%%%%%%%%%%%%%%%%%%%%%%%%%%%%%%%%%%%%%%%%%%%%%%%%%%%%%
\subsection{Common False Reciprocal Identity}
\label{subsec:false-reciprocal-sum}
%%%%%%%%%%%%%%%%%%%%%%%%%%%%%%%%%%%%%%%%%%%%%%%%%%%%%%%%%%%%

In general,

\[
\boxed{
\frac1{a+b}
\neq
\frac1a+\frac1b.
}
\]

Instead,

\[
\boxed{
\frac1a+\frac1b
=
\frac{a+b}{ab}.
}
\]

%%%%%%%%%%%%%%%%%%%%%%%%%%%%%%%%%%%%%%%%%%%%%%%%%%%%%%%%%%%%
\subsection{Common False Cancellation With Powers}
\label{subsec:false-power-cancellation}
%%%%%%%%%%%%%%%%%%%%%%%%%%%%%%%%%%%%%%%%%%%%%%%%%%%%%%%%%%%%

In general,

\[
\boxed{
\frac{x^2+1}{x}
\neq
x+1.
}
\]

The correct simplification is

\[
\boxed{
\frac{x^2+1}{x}
=
x+\frac1x,
\qquad
x\neq0.
}
\]

%%%%%%%%%%%%%%%%%%%%%%%%%%%%%%%%%%%%%%%%%%%%%%%%%%%%%%%%%%%%
\subsection{Common False Square-Root Cancellation}
\label{subsec:false-square-root-cancellation}
%%%%%%%%%%%%%%%%%%%%%%%%%%%%%%%%%%%%%%%%%%%%%%%%%%%%%%%%%%%%

From

\[
\sqrt{x^2},
\]

the correct simplification over \(\mathbb{R}\) is

\[
\boxed{
|x|,
}
\]

not necessarily \(x\).

%%%%%%%%%%%%%%%%%%%%%%%%%%%%%%%%%%%%%%%%%%%%%%%%%%%%%%%%%%%%
\subsection{Algebraic Identities in Calculus}
\label{subsec:algebraic-identities-calculus}
%%%%%%%%%%%%%%%%%%%%%%%%%%%%%%%%%%%%%%%%%%%%%%%%%%%%%%%%%%%%

Algebraic identities frequently simplify limits and derivatives.

For example,

\[
\frac{x^2-a^2}{x-a}
=
x+a
\]

for \(x\neq a\).

Thus

\[
\boxed{
\lim_{x\to a}
\frac{x^2-a^2}{x-a}
=
2a.
}
\]

%%%%%%%%%%%%%%%%%%%%%%%%%%%%%%%%%%%%%%%%%%%%%%%%%%%%%%%%%%%%
\subsection{Algebraic Identities in Integration}
\label{subsec:algebraic-identities-integration}
%%%%%%%%%%%%%%%%%%%%%%%%%%%%%%%%%%%%%%%%%%%%%%%%%%%%%%%%%%%%

Factoring and partial fractions are central to integration.

For example,

\[
\boxed{
\frac{1}{x^2-1}
=
\frac12
\left(
\frac1{x-1}
-
\frac1{x+1}
\right).
}
\]

This transforms a rational integrand into simpler pieces.

%%%%%%%%%%%%%%%%%%%%%%%%%%%%%%%%%%%%%%%%%%%%%%%%%%%%%%%%%%%%
\subsection{Algebraic Identities in Number Theory}
\label{subsec:algebraic-identities-number-theory}
%%%%%%%%%%%%%%%%%%%%%%%%%%%%%%%%%%%%%%%%%%%%%%%%%%%%%%%%%%%%

The identity

\[
\boxed{
a^n-b^n
=
(a-b)
\left(
a^{n-1}+\cdots+b^{n-1}
\right)
}
\]

implies divisibility facts such as

\[
\boxed{
a-b
\mid
a^n-b^n.
}
\]

Thus polynomial identities produce arithmetic consequences.

%%%%%%%%%%%%%%%%%%%%%%%%%%%%%%%%%%%%%%%%%%%%%%%%%%%%%%%%%%%%
\subsection{Algebraic Identities in Probability}
\label{subsec:algebraic-identities-probability}
%%%%%%%%%%%%%%%%%%%%%%%%%%%%%%%%%%%%%%%%%%%%%%%%%%%%%%%%%%%%

Variance calculations use

\[
\boxed{
(X-\mu)^2
=
X^2-2\mu X+\mu^2.
}
\]

Taking expectations gives

\[
\boxed{
\operatorname{Var}(X)
=
\mathbb{E}[X^2]
-
(\mathbb{E}[X])^2.
}
\]

%%%%%%%%%%%%%%%%%%%%%%%%%%%%%%%%%%%%%%%%%%%%%%%%%%%%%%%%%%%%
\subsection{Algebraic Identities in Linear Algebra}
\label{subsec:algebraic-identities-linear-algebra}
%%%%%%%%%%%%%%%%%%%%%%%%%%%%%%%%%%%%%%%%%%%%%%%%%%%%%%%%%%%%

Matrix identities require attention to multiplication order.

For example,

\[
\boxed{
(AB)^{-1}
=
B^{-1}A^{-1},
}
\]

not

\[
A^{-1}B^{-1}
\]

in general.

Similarly,

\[
\boxed{
(AB)^T
=
B^TA^T.
}
\]

%%%%%%%%%%%%%%%%%%%%%%%%%%%%%%%%%%%%%%%%%%%%%%%%%%%%%%%%%%%%
\subsection{Algebraic Identities in Optimization}
\label{subsec:algebraic-identities-optimization}
%%%%%%%%%%%%%%%%%%%%%%%%%%%%%%%%%%%%%%%%%%%%%%%%%%%%%%%%%%%%

Quadratic objectives are often rewritten by completing the square.

For example,

\[
x^2-6x+11
\]

becomes

\[
\boxed{
(x-3)^2+2.
}
\]

The minimum is therefore immediately visible:

\[
\boxed{
\min_x
(x^2-6x+11)
=
2
}
\]

at

\[
\boxed{
x=3.
}
\]

%%%%%%%%%%%%%%%%%%%%%%%%%%%%%%%%%%%%%%%%%%%%%%%%%%%%%%%%%%%%
\subsection{Algebraic Identities in Complex Analysis}
\label{subsec:algebraic-identities-complex-analysis}
%%%%%%%%%%%%%%%%%%%%%%%%%%%%%%%%%%%%%%%%%%%%%%%%%%%%%%%%%%%%

Complex conjugation gives

\[
\boxed{
\frac{1}{z}
=
\frac{\overline z}{|z|^2},
\qquad
z\neq0.
}
\]

This follows from

\[
z\overline z
=
|z|^2.
\]

%%%%%%%%%%%%%%%%%%%%%%%%%%%%%%%%%%%%%%%%%%%%%%%%%%%%%%%%%%%%
\subsection{Algebraic Identities Workflow}
\label{subsec:algebraic-identities-workflow}
%%%%%%%%%%%%%%%%%%%%%%%%%%%%%%%%%%%%%%%%%%%%%%%%%%%%%%%%%%%%

When simplifying an expression:

\begin{enumerate}
    \item identify the domain;
    \item look for common factors;
    \item identify powers or radicals;
    \item recognize special polynomial patterns;
    \item factor before canceling;
    \item simplify rational expressions carefully;
    \item use conjugates when radicals appear in denominators;
    \item apply logarithm rules only on valid positive arguments;
    \item preserve excluded values;
    \item avoid replacing identities with false pattern matching;
    \item verify the final expression numerically when appropriate.
\end{enumerate}

%%%%%%%%%%%%%%%%%%%%%%%%%%%%%%%%%%%%%%%%%%%%%%%%%%%%%%%%%%%%
\subsection{Quick Algebraic Identity Table}
\label{subsec:quick-algebraic-identity-table}
%%%%%%%%%%%%%%%%%%%%%%%%%%%%%%%%%%%%%%%%%%%%%%%%%%%%%%%%%%%%

\[
\boxed{
(a+b)^2
=
a^2+2ab+b^2
}
\]

\[
\boxed{
(a-b)^2
=
a^2-2ab+b^2
}
\]

\[
\boxed{
a^2-b^2
=
(a-b)(a+b)
}
\]

\[
\boxed{
(a+b)^3
=
a^3+3a^2b+3ab^2+b^3
}
\]

\[
\boxed{
(a-b)^3
=
a^3-3a^2b+3ab^2-b^3
}
\]

\[
\boxed{
a^3+b^3
=
(a+b)(a^2-ab+b^2)
}
\]

\[
\boxed{
a^3-b^3
=
(a-b)(a^2+ab+b^2)
}
\]

\[
\boxed{
a^m a^n
=
a^{m+n}
}
\]

\[
\boxed{
\frac{a^m}{a^n}
=
a^{m-n}
}
\]

\[
\boxed{
(a^m)^n
=
a^{mn}
}
\]

\[
\boxed{
(ab)^n
=
a^nb^n
}
\]

\[
\boxed{
\log_b(xy)
=
\log_bx+\log_by
}
\]

\[
\boxed{
\log_b(x/y)
=
\log_bx-\log_by
}
\]

\[
\boxed{
\log_b(x^r)
=
r\log_bx
}
\]

\[
\boxed{
\sum_{k=1}^{n}k
=
\frac{n(n+1)}{2}
}
\]

\[
\boxed{
\sum_{k=1}^{n}k^2
=
\frac{n(n+1)(2n+1)}{6}
}
\]

\[
\boxed{
\sum_{k=1}^{n}k^3
=
\left(
\frac{n(n+1)}{2}
\right)^2
}
\]

\[
\boxed{
\sum_{k=0}^{n-1}ar^k
=
a\frac{1-r^n}{1-r},
\qquad
r\neq1
}
\]

\[
\boxed{
\sum_{k=0}^{\infty}ar^k
=
\frac{a}{1-r},
\qquad
|r|<1.
}
\]

%%%%%%%%%%%%%%%%%%%%%%%%%%%%%%%%%%%%%%%%%%%%%%%%%%%%%%%%%%%%
\subsection{Common False Identity Table}
\label{subsec:common-false-identity-table}
%%%%%%%%%%%%%%%%%%%%%%%%%%%%%%%%%%%%%%%%%%%%%%%%%%%%%%%%%%%%

The following are generally false:

\[
\boxed{
(a+b)^2
\neq
a^2+b^2,
}
\]

\[
\boxed{
\sqrt{a+b}
\neq
\sqrt a+\sqrt b,
}
\]

\[
\boxed{
\frac1{a+b}
\neq
\frac1a+\frac1b,
}
\]

\[
\boxed{
\ln(a+b)
\neq
\ln a+\ln b,
}
\]

\[
\boxed{
|a+b|
\neq
|a|+|b|,
}
\]

and

\[
\boxed{
(AB)^{-1}
\neq
A^{-1}B^{-1}
}
\]

in general.

%%%%%%%%%%%%%%%%%%%%%%%%%%%%%%%%%%%%%%%%%%%%%%%%%%%%%%%%%%%%
\subsection{Exercises}
\label{subsec:algebraic-identities-exercises}
%%%%%%%%%%%%%%%%%%%%%%%%%%%%%%%%%%%%%%%%%%%%%%%%%%%%%%%%%%%%

\begin{exercise}[1]
Expand

\[
(x+4)^2.
\]
\end{exercise}

\begin{exercise}[1]
Expand

\[
(2x-3)^2.
\]
\end{exercise}

\begin{exercise}[1]
Factor

\[
x^2-49.
\]
\end{exercise}

\begin{exercise}[1]
Factor

\[
x^3-27.
\]
\end{exercise}

\begin{exercise}[1]
Factor

\[
x^3+64.
\]
\end{exercise}

\begin{exercise}[2]
Expand

\[
(2x+y)^3.
\]
\end{exercise}

\begin{exercise}[2]
Factor

\[
16x^4-81.
\]
\end{exercise}

\begin{exercise}[2]
Factor completely

\[
x^4-16.
\]
\end{exercise}

\begin{exercise}[2]
Simplify

\[
x^3x^{-5}.
\]
\end{exercise}

\begin{exercise}[2]
Simplify

\[
\frac{x^8}{x^3}.
\]
\end{exercise}

\begin{exercise}[2]
Simplify

\[
(x^3)^4.
\]
\end{exercise}

\begin{exercise}[2]
Rewrite

\[
x^{-4}
\]

using positive exponents.
\end{exercise}

\begin{exercise}[2]
Simplify

\[
\sqrt{72}.
\]
\end{exercise}

\begin{exercise}[2]
Simplify

\[
\sqrt{x^2}
\]

over the real numbers.
\end{exercise}

\begin{exercise}[2]
Rationalize

\[
\frac{1}{\sqrt5}.
\]
\end{exercise}

\begin{exercise}[2]
Rationalize

\[
\frac{1}{2+\sqrt3}.
\]
\end{exercise}

\begin{exercise}[2]
Expand

\[
\ln(x^3y^2)
\]

assuming \(x,y>0\).
\end{exercise}

\begin{exercise}[2]
Condense

\[
2\ln x-\ln y
\]

into one logarithm.
\end{exercise}

\begin{exercise}[2]
Use change of base to rewrite

\[
\log_5x
\]

using natural logarithms.
\end{exercise}

\begin{exercise}[2]
Explain why

\[
\ln(x+y)
=
\ln x+\ln y
\]

is generally false.
\end{exercise}

\begin{exercise}[2]
Complete the square:

\[
x^2+8x+3.
\]
\end{exercise}

\begin{exercise}[2]
Complete the square:

\[
2x^2-12x+5.
\]
\end{exercise}

\begin{exercise}[2]
Solve

\[
3x^2-5x-2=0
\]

using the quadratic formula.
\end{exercise}

\begin{exercise}[2]
Use Vieta's formulas to find the sum and product of roots of

\[
2x^2-7x+3.
\]
\end{exercise}

\begin{exercise}[2]
Expand

\[
(x+y)^5
\]

using the binomial theorem.
\end{exercise}

\begin{exercise}[2]
Compute

\[
\binom{8}{3}.
\]
\end{exercise}

\begin{exercise}[3]
Prove Pascal's identity.
\end{exercise}

\begin{exercise}[3]
Prove

\[
\sum_{k=0}^{n}
\binom{n}{k}
=
2^n.
\]
\end{exercise}

\begin{exercise}[3]
Prove

\[
\sum_{k=0}^{n}
(-1)^k
\binom{n}{k}
=
0
\]

for \(n\geq1\).
\end{exercise}

\begin{exercise}[2]
Simplify

\[
\frac{x^2-9}{x-3}
\]

and state the excluded value.
\end{exercise}

\begin{exercise}[2]
Simplify

\[
\frac{x^2+x}{x}.
\]
\end{exercise}

\begin{exercise}[2]
Simplify

\[
\frac{
\frac1x+\frac1y
}{
\frac1x-\frac1y
}.
\]
\end{exercise}

\begin{exercise}[2]
Show that

\[
|ab|
=
|a||b|.
\]
\end{exercise}

\begin{exercise}[3]
Prove the reverse triangle inequality.
\end{exercise}

\begin{exercise}[2]
Compute

\[
\overline{(3+4i)(2-i)}
\]

using conjugate identities.
\end{exercise}

\begin{exercise}[2]
Show that

\[
\operatorname{Re}(z)
=
\frac{z+\overline z}{2}.
\]
\end{exercise}

\begin{exercise}[2]
Show that

\[
\operatorname{Im}(z)
=
\frac{z-\overline z}{2i}.
\]
\end{exercise}

\begin{exercise}[2]
Derive

\[
\cos\theta
=
\frac{e^{i\theta}+e^{-i\theta}}{2}.
\]
\end{exercise}

\begin{exercise}[2]
Derive

\[
\sin\theta
=
\frac{e^{i\theta}-e^{-i\theta}}{2i}.
\]
\end{exercise}

\begin{exercise}[2]
Verify

\[
\cosh^2x-\sinh^2x=1.
\]
\end{exercise}

\begin{exercise}[2]
Derive

\[
x^2+\frac1{x^2}
=
\left(
x+\frac1x
\right)^2-2.
\]
\end{exercise}

\begin{exercise}[3]
Express

\[
x^3+\frac1{x^3}
\]

in terms of

\[
x+\frac1x.
\]
\end{exercise}

\begin{exercise}[2]
Evaluate

\[
1+2+\cdots+100.
\]
\end{exercise}

\begin{exercise}[2]
Evaluate

\[
1^2+2^2+\cdots+10^2.
\]
\end{exercise}

\begin{exercise}[2]
Evaluate

\[
1^3+2^3+\cdots+10^3.
\]
\end{exercise}

\begin{exercise}[2]
Evaluate

\[
1+3+9+\cdots+3^8.
\]
\end{exercise}

\begin{exercise}[2]
Evaluate

\[
\sum_{k=0}^{\infty}
\left(
\frac12
\right)^k.
\]
\end{exercise}

\begin{exercise}[3]
Show that

\[
\frac{1}{k(k+1)}
=
\frac1k-\frac1{k+1}.
\]
\end{exercise}

\begin{exercise}[3]
Use the previous identity to evaluate

\[
\sum_{k=1}^{n}
\frac{1}{k(k+1)}.
\]
\end{exercise}

\begin{exercise}[3]
Simplify

\[
\frac{x^5-a^5}{x-a}.
\]
\end{exercise}

\begin{exercise}[3]
Derive the difference quotient of \(x^4\) from the binomial theorem.
\end{exercise}

\begin{exercise}[3]
Show

\[
(a+b)^2+(a-b)^2
=
2a^2+2b^2.
\]
\end{exercise}

\begin{exercise}[3]
Prove the parallelogram identity in \(\R^n\).
\end{exercise}

\begin{exercise}[3]
Verify

\[
(AB)^T
=
B^TA^T.
\]
\end{exercise}

\begin{exercise}[3]
Verify

\[
(AB)^{-1}
=
B^{-1}A^{-1}
\]

for invertible matrices.
\end{exercise}

\begin{exercise}[3]
Show that

\[
\operatorname{tr}(AB)
=
\operatorname{tr}(BA).
\]
\end{exercise}

\begin{exercise}[3]
Verify the finite matrix geometric-series identity

\[
(I-A)
\sum_{k=0}^{n-1}A^k
=
I-A^n.
\]
\end{exercise}

\begin{exercise}[3]
Explain why

\[
\frac{x^2-1}{x-1}
\]

and \(x+1\) are not identical as functions on their natural domains.
\end{exercise}

\begin{exercise}[3]
Give an example showing how squaring both sides of an equation creates an
extraneous solution.
\end{exercise}

\begin{exercise}[3]
Give an example showing why dividing by a variable can lose a valid solution.
\end{exercise}

\begin{exercise}[3]
Explain why

\[
(a+b)^n=a^n+b^n
\]

is generally false.
\end{exercise}

\begin{exercise}[3]
Explain why

\[
\sqrt{a+b}
=
\sqrt a+\sqrt b
\]

is generally false.
\end{exercise}

\begin{exercise}[3]
Explain why

\[
\frac1{a+b}
=
\frac1a+\frac1b
\]

is generally false.
\end{exercise}

\begin{exercise}[4]
Prove the binomial theorem by induction.
\end{exercise}

\begin{exercise}[4]
Derive the multinomial theorem from repeated use of the binomial theorem.
\end{exercise}

\begin{exercise}[4]
Prove the general factorization of \(a^n-b^n\).
\end{exercise}

\begin{exercise}[4]
Prove the factorization of \(a^n+b^n\) for odd \(n\).
\end{exercise}

\begin{exercise}[4]
Derive the quadratic formula by completing the square.
\end{exercise}

\begin{exercise}[4]
Prove Vieta's formulas for a quadratic polynomial.
\end{exercise}

\begin{exercise}[4]
Prove the formula for the sum of the first \(n\) squares.
\end{exercise}

\begin{exercise}[4]
Prove the formula for the sum of the first \(n\) cubes.
\end{exercise}

\begin{exercise}[4]
Derive the finite geometric-series formula algebraically.
\end{exercise}

\begin{exercise}[4]
Prove the infinite geometric-series formula using limits.
\end{exercise}

\begin{exercise}[4]
Derive the Neumann series for a matrix under a suitable norm condition.
\end{exercise}

\begin{exercise}[4]
Study the role of algebraic identities in rational-function integration.
\end{exercise}

\begin{exercise}[4]
Study how difference-of-powers factorization is used in number-theoretic
divisibility proofs.
\end{exercise}

%%%%%%%%%%%%%%%%%%%%%%%%%%%%%%%%%%%%%%%%%%%%%%%%%%%%%%%%%%%%
\subsection{Conceptual Summary}
\label{subsec:algebraic-identities-summary}
%%%%%%%%%%%%%%%%%%%%%%%%%%%%%%%%%%%%%%%%%%%%%%%%%%%%%%%%%%%%

Algebraic identities provide reusable symbolic relationships that simplify
expressions and expose mathematical structure.

Fundamental examples include

\[
\boxed{
(a+b)^2
=
a^2+2ab+b^2,
}
\]

\[
\boxed{
a^2-b^2
=
(a-b)(a+b),
}
\]

\[
\boxed{
a^3-b^3
=
(a-b)(a^2+ab+b^2),
}
\]

and

\[
\boxed{
(a+b)^n
=
\sum_{k=0}^{n}
\binom{n}{k}
a^{n-k}b^k.
}
\]

Exponent rules give

\[
\boxed{
a^ma^n=a^{m+n},
}
\]

while logarithms convert multiplication into addition:

\[
\boxed{
\log_b(xy)
=
\log_bx+\log_by.
}
\]

Rational-expression identities require attention to nonzero denominators,
and radical identities require attention to domain and principal-root
conventions.

Complex-number identities include

\[
\boxed{
|z|^2
=
z\overline z.
}
\]

Linear algebra introduces noncommutative identities such as

\[
\boxed{
(AB)^{-1}
=
B^{-1}A^{-1}.
}
\]

The most important practical principle is

\[
\boxed{
\text{simplify algebraically without changing the domain or logical meaning}.
}
\]

%%%%%%%%%%%%%%%%%%%%%%%%%%%%%%%%%%%%%%%%%%%%%%%%%%%%%%%%%%%%
\subsection{Section Connections}
\label{subsec:algebraic-identities-connections}
%%%%%%%%%%%%%%%%%%%%%%%%%%%%%%%%%%%%%%%%%%%%%%%%%%%%%%%%%%%%

\chapterconnection{
Algebraic Identities collects formulas that appear throughout nearly every
chapter of the textbook. Factoring identities support Algebra and Number
Theory, difference quotients support Calculus, binomial identities connect
Algebra with Combinatorics and Probability, logarithmic and exponential
identities support Analysis and Differential Equations, conjugate identities
support Complex Analysis, and matrix identities support Linear Algebra,
Optimization, Numerical Mathematics, and Applied Modeling. The next reference
section, Geometric Formulas, collects perimeter, area, surface-area, volume,
distance, coordinate, angle, and related formulas for common geometric
objects.
}

%%%%%%%%%%%%%%%%%%%%%%%%%%%%%%%%%%%%%%%%%%%%%%%%%%%%%%%%%%%%
% SECTION SOURCE TRACKING
%%%%%%%%%%%%%%%%%%%%%%%%%%%%%%%%%%%%%%%%%%%%%%%%%%%%%%%%%%%%

%------------------------------------------------------------
% 14.5 Algebraic Identities
%------------------------------------------------------------
%
% Source basis:
%   - Standard arithmetic and algebraic identities.
%   - Standard exponent, radical, logarithm, polynomial,
%     combinatorial, rational-expression, complex-number,
%     sequence, and matrix identities.
%
% Used for:
%   - commutative laws;
%   - associative laws;
%   - distributive law;
%   - negation rules;
%   - exponent laws;
%   - rational exponents;
%   - radical identities;
%   - rationalization;
%   - logarithm rules;
%   - squares and cubes;
%   - difference of squares;
%   - sums and differences of cubes;
%   - general differences of powers;
%   - perfect-square trinomials;
%   - completing the square;
%   - quadratic formula;
%   - discriminants;
%   - Vieta formulas;
%   - Remainder and Factor Theorems;
%   - binomial theorem;
%   - Pascal identities;
%   - multinomial theorem;
%   - rational expressions;
%   - absolute-value identities;
%   - triangle inequalities;
%   - complex conjugates;
%   - Euler formulas;
%   - hyperbolic identities;
%   - reciprocal identities;
%   - arithmetic and geometric sums;
%   - sums of powers;
%   - telescoping identities;
%   - partial fractions;
%   - difference quotients;
%   - norm and inner-product identities;
%   - matrix transpose, inverse, trace, and determinant identities;
%   - Neumann series;
%   - domain restrictions;
%   - common false identities;
%   - applications and exercises.
%
% Notes:
%   - Geometric Formulas is developed next in Section 14.6.
%   - Trigonometric Identities follows in Section 14.7.
%
%%%%%%%%%%%%%%%%%%%%%%%%%%%%%%%%%%%%%%%%%%%%%%%%%%%%%%%%%%%%